%%% Fundamental｜基础知识 %%%


\section*{2020 年}


\textbf{选择题｜3 · 3 分}

设 $f(x) = \ln(x + \sqrt{1 + x^2})$，\\
① $f(x)$ 是奇函数； \quad ② $f(x)$ 的反函数是 $y = \frac{e^x-e^{-x}}{2}$；\\
③ $f'(x) = \frac{1}{\sqrt{1+x^2}}$； \quad ④ $f(x)$ 的一个原函数是 $x \ln(x + \sqrt{1 + x^2}) - \sqrt{1 + x^2}$。\\
则以上叙述不正确的个数是( \quad )。\\
A. 0 个； \quad B. 1 个； \quad C. 2 个； \quad D. 3 个。


\textbf{填空题｜4 · 3 分}

曲线 $\begin{cases} x^2 + 3z^2 = 1 \\ y = 0 \end{cases}$ 绕 $z$ 轴旋转一周所得到的旋转曲面方程为 \underline{\hspace{3cm}}。


\textbf{计算题｜7 · 8 分}

求过直线 $l: \frac{x-2}{1} = \frac{y-1}{2} = \frac{z-3}{3}$ 和点 $M(2,0,1)$ 的平面的方程。


\section*{2021 年}


\textbf{选择题｜6 · 3 分}

过点 $A(1,1, 2)$，$B(2,1,10)$，$C(0,2,3)$ 的平面方程为（ \quad ）。\\
A. $8x + 9y - z - 15 = 0$ \quad B. $8x - 9y + z - 15 = 0$\\
C. $8x + 9y + z - 15 = 0$ \quad D. $8x - 9y - z - 15 = 0$


\textbf{选择题｜7 · 3 分}

设向量 $\vec{a}, \vec{b}$ 为非零向量且 $\vec{a} \perp \vec{b}$，则必有（ \quad ）。\\
A. $|\vec{a} + \vec{b}| = |\vec{a}| + |\vec{b}|$ \quad B. $|\vec{a} - \vec{b}| = |\vec{a}| - |\vec{b}|$\\
C. $|\vec{a} - \vec{b}| = |\vec{a} + \vec{b}|$ \quad D. $\vec{a} - \vec{b} = \vec{a} + \vec{b}$


\section*{2022 年}


\textbf{选择题｜6 · 3 分}

设空间三点的坐标分别为 $A(1, 1, 1)$，$B(2, 2, 1)$，$C(2, 1, 2)$，则 $\angle BAC =$（ \quad ）\\
A. $\pi$ \quad B. $\frac{\pi}{4}$ \quad C. $\frac{\pi}{2}$ \quad D. $\frac{\pi}{3}$


\section*{2023 年}


\textbf{选择题｜3 · 3 分}

设 $f(x)$ 在 $(-\infty, +\infty)$ 内是奇函数且可导，则下列函数是奇函数的是（ \quad ）。\\
A. $\int_0^x (\sin x)f(t) dt$ \quad B. $\int_0^x [\sin t + f(t)] dt$ \quad C. $\int_0^x f(\sin t) dt$ \quad D. $\sin[f'(x)]$


\textbf{填空题｜5 · 3 分}

过点 $(3, 1, -2)$ 且过直线 $\frac{x-4}{0} = \frac{y+3}{2} = z$ 的平面方程为 \underline{\hspace{3cm}}。


\section*{2024 年}


\textbf{选择题｜5 · 3 分}

下列曲线中，绕 $y$ 轴旋转而成椭球面 $3x^2 + 2y^2 + 3z^2 = 1$ 的曲线是（ \quad ）\\
A. $\begin{cases} 2x^2 + 3y^2 = 1 \\ z = 0 \end{cases}$ \quad B. $\begin{cases} 3y^2 + 2z^2 = 1 \\ x = 0 \end{cases}$\\
C. $\begin{cases} 3x^2 + 2y^2 = 1 \\ z = 0 \end{cases}$ \quad D. $\begin{cases} 3x^2 + 3z^2 = 1 \\ y = 0 \end{cases}$


\textbf{计算题｜6 · 9 分}

求通过点 $(1, 2, -1)$ 且与直线 $\begin{cases} 2x - 3y + z - 5 = 0 \\ 3x + y - 2z - 4 = 0 \end{cases}$ 垂直的平面方程。
